\chapter{Conclusiones y recomendaciones}
\label{cap:conclusiones}

Empezamos con una pregunta práctica: ¿cómo medir un pez a través del vidrio de una pecera con una sola cámara? Para responderla pusimos a prueba, en un montaje controlado aire--vidrio--agua, cuatro rutas de conversión métrica sobre las mismas imágenes y los mismos puntos anatómicos, y medimos cuánto se ensancha la incertidumbre cuando no conocemos la profundidad. Trabajamos con seis réplicas rígidas de PLA, cuatro profundidades y tres capturas técnicas por condición: 72 imágenes en total. Todo lo que sigue se refiere a ese conjunto y a las referencias físicas registradas en el CSV maestro.

El aporte se puede condensar en una regla de decisión (Tabla~\ref{tab:regla} y Figura~\ref{fig:5-1}): según lo que se sepa de la escena, qué conversión conviene usar y cuánto error esperar en este montaje.

\begin{table}[htbp]
\centering\small
\caption{Regla de decisión: qué conversión usar según la situación, con el error observado en este montaje.}
\label{tab:regla}
\begin{tabularx}{\textwidth}{@{}X p{2.6cm} p{3.0cm} X@{}}
\toprule
Situación de uso & Método recomendado & Error esperado en este montaje & Condición \\
\midrule
Se puede poner un marcador junto al pez, en su mismo plano & \Aprime{} (homografía con marcador coplanar) & $\approx$ 2--3~\% de LT (3,3~mm) & Coplanaridad dentro de $\pm 1$~cm; no requiere calibración \\
No se puede poner un marcador junto al pez, pero se conoce su distancia al vidrio (riel, guía, tanque angosto) & C (puerto plano) & $< 1$~\% de LT (1,2~mm) & Calibración en aire y parámetros $d$, $t_g$ e índices del puerto \\
Hay calibración bajo agua en el mismo tanque y se conoce la posición del pez en el marco de esa calibración (riel o guía) & B (pinhole efectivo) & $< 1$~\% de LT (1,1~mm) & Damero sumergido y referencia del plano; comprobado hasta 18~cm de la zona calibrada \\
Marcador fijo en el vidrio y pez a distancia desconocida & Ninguno es aceptable (A) & 10--40~\% según la distancia & Solo sirve si el pez está pegado al vidrio \\
No se sabe nada de la profundidad & C con intervalo $\IL$ & 56--82~\% de ancho: no utilizable & Hay que acotar $z$; con $\pm 1$~cm el ancho baja a 4--5~\% \\
\bottomrule
\end{tabularx}
\end{table}

\begin{figure}[htbp]
\centering
\begin{tikzpicture}[>=Stealth,font=\small,
  preg/.style={draw,rounded corners=3pt,fill=gray!8,text width=5.2cm,align=center,minimum height=9mm},
  hoja/.style={draw,rounded corners=3pt,text width=4.6cm,align=center,minimum height=9mm,font=\small\bfseries}]
  \node[preg] (q1) at (0,0) {¿Se puede colocar un marcador coplanar con el pez?};
  \node[hoja,fill=metAp!15] (h1) at (6.3,0) {\Aprime\\ \normalfont error $\approx$ 2--3~\%};
  \node[preg] (q2) at (0,-2.2) {¿Se conoce la distancia pez--vidrio con $\pm 1$~cm?};
  \node[hoja,fill=metC!15] (h2) at (6.3,-2.2) {C, o B con calibración bajo agua\\ \normalfont error $< 1$~\%};
  \node[preg] (q3) at (0,-4.4) {¿Se puede acotar $z$ con una guía o una segunda vista?};
  \node[hoja,fill=metC!15] (h3) at (6.3,-4.4) {C con intervalo $\IL$\\ \normalfont ancho 4--17~\% según la cota};
  \node[hoja,fill=metA!15] (h4) at (0,-6.6) {No medir\\ \normalfont intervalo de 56--82~\%: no utilizable};
  \draw[->] (q1) -- node[above,font=\footnotesize]{sí} (h1);
  \draw[->] (q1) -- node[right,font=\footnotesize]{no} (q2);
  \draw[->] (q2) -- node[above,font=\footnotesize]{sí} (h2);
  \draw[->] (q2) -- node[right,font=\footnotesize]{no} (q3);
  \draw[->] (q3) -- node[above,font=\footnotesize]{sí} (h3);
  \draw[->] (q3) -- node[right,font=\footnotesize]{no} (h4);
\end{tikzpicture}
\caption[Diagrama de decisión para la conversión métrica monocular a través de un vidrio plano]{Diagrama de decisión para la conversión métrica monocular a través de un vidrio plano. Cada hoja indica el error observado en este montaje (Capítulo~\ref{cap:resultados}); B requiere una calibración bajo agua en el mismo tanque. Fuente: elaboración propia.}
\label{fig:5-1}
\end{figure}

\section{Conclusiones por objetivo}
\label{sec:conclusiones}

\begin{enumerate}
  \item \textbf{Conjunto de datos (objetivo 1).} Construimos y documentamos un conjunto de 72 imágenes de seis réplicas a cuatro profundidades, con referencias físicas de calibre, una tabla maestra que permite rastrear cada valor por réplica, profundidad, repetición e imagen, y un registro de exclusiones. Obtuvimos valores numéricos para 284 de los 288 registros; los cuatro restantes corresponden a A y \Aprime{} en dos imágenes en las que no se detectó el marcador. Resumimos las repeticiones técnicas por combinación réplica--profundidad. Objetivo cumplido.
  \item \textbf{Protocolo y repetibilidad del marcado (objetivo 2).} Definimos y aplicamos un protocolo de marcado manual único por imagen, con cuatro puntos anatómicos que comparten las cuatro rutas. No cuantificamos en esta versión la repetibilidad del marcado, que queda como verificación inmediata. Objetivo cumplido parcialmente.
  \item \textbf{Implementación y validación de las rutas (objetivo 3).} Implementamos las cuatro rutas en una aplicación propia de código abierto, con pruebas unitarias de sanidad para el módulo de puerto plano, y las evaluamos contra las referencias de calibre de las seis réplicas en 72 imágenes. B logró un MAE de 1,08~mm (0,73~\%) en LT y de 0,38~mm (1,26~\%) en PC; C, de 1,22~mm (0,88~\%) y 0,56~mm (1,62~\%); y \Aprime, de 3,29~mm (2,28~\%) y 1,03~mm (2,78~\%). No hicimos la comprobación instrumental con barra patrón. Objetivo cumplido en cuanto a la validación contra referencias físicas; la comprobación con patrón queda como trabajo futuro.
  \item \textbf{Comparación pareada según la separación (objetivo 4).} A fue la ruta más sensible a la profundidad: su MAE de LT pasó de 13,77~mm a 5~cm a 57,64~mm a 23~cm (2,43~mm por centímetro), tal como predice la geometría del puerto plano. B y C se mantuvieron entre 1,0 y 1,4~mm en las cuatro posiciones, sin tendencia con la separación. C superó a A y \Aprime{} en 6 de 6 réplicas ($p = 0{,}031$) y en 24 de 24 combinaciones, pero no se diferenció de B (3 de 6 réplicas, $p = 1$). H1 se sostiene frente a A y \Aprime{} y no frente a B. H2 no se sostiene: hasta 18~cm de la zona de calibración, B no se degradó con la separación. Objetivo cumplido.
  \item \textbf{Intervalo por profundidad desconocida (objetivo 5).} Calculamos el intervalo $\IL$ con C sobre las 72 capturas, y coincidió con la predicción geométrica del montaje (razón media 1,006). Su ancho es de 56--82~\% de LT si $z$ es desconocida en toda la pecera, de 11--16~\% con $\pm 3$~cm y de 4--5~\% con $\pm 1$~cm. La cobertura empírica fue de 1,00 en toda la pecera y con $\pm 3$~cm, y de 0,81 con $\pm 1$~cm. H3 se sostiene. Objetivo cumplido.
  \item \textbf{Sensibilidad (objetivo 6).} Cuantificamos cuánto depende C de la distancia cámara--vidrio y del índice del agua (Tabla~\ref{tab:sensibilidad}). Un error de $\pm 20$~\% en $d$ cambió la LT en 14~\% (cerca de 0,3~\% por milímetro); un cambio de $\pm 0{,}01$ en $\mu_w$, apenas 0,23~\%. No cuantificamos en esta versión la sensibilidad a la repetibilidad del marcado. Objetivo cumplido parcialmente.
\end{enumerate}

La Tabla~\ref{tab:cumplimiento} resume el grado de cumplimiento de cada objetivo.

\begin{table}[htbp]
\centering\small
\caption{Matriz de cumplimiento de los objetivos específicos.}
\label{tab:cumplimiento}
\begin{tabularx}{\textwidth}{@{}p{1.6cm} X p{1.8cm} X@{}}
\toprule
Objetivo & Resultado que lo sostiene & Grado & Pendiente \\
\midrule
1 & Tablas~\ref{tab:replicas} y~\ref{tab:registros}; repositorio & Cumplido & Publicar los datos con la etiqueta \texttt{\RepoTag} \\
2 & Sección~\ref{sec:procedimiento} & Parcial & Remarcado de doce imágenes \\
3 & Tablas~\ref{tab:calibraciones} y~\ref{tab:agregados} & Cumplido & Comprobación con barra patrón \\
4 & Tabla~\ref{tab:posiciones-mae}; Tabla~\ref{tab:hipotesis} & Cumplido & --- \\
5 & Tabla~\ref{tab:intervalo}; Figuras~\ref{fig:4-3} y~\ref{fig:4-Lz} & Cumplido & --- \\
6 & Tabla~\ref{tab:sensibilidad} & Parcial & Efecto del remarcado sobre las medidas \\
\bottomrule
\end{tabularx}
\end{table}

En pocas palabras: a través de un vidrio plano, la escala de un marcador pegado al vidrio no sirve para medir un pez que nada a otra profundidad, porque el error crece 2,4~mm por cada centímetro de separación. Un marcador en el plano del pez baja ese error a 2--3~\%. Un modelo de cámara calibrado, con la profundidad conocida, lo baja a menos del 1~\%, y da lo mismo si es el pinhole efectivo calibrado bajo el agua (B) o el puerto plano (C). Cuando la profundidad no se conoce, la geometría encierra la longitud en un intervalo demasiado ancho para hacer morfometría, salvo que la profundidad se acote a $\pm 1$~cm. En un sistema de una sola cámara, entonces, el problema a resolver es la profundidad, no la refracción, y este trabajo pone un número a cuánto pesa.

\section{Limitaciones}

Trabajamos con réplicas rígidas, una sola cámara y una pecera de vidrio, con la iluminación y el encuadre bajo control. No evaluamos peces vivos, movimiento, agua turbia, oclusiones ni posiciones laterales planificadas. Seis réplicas en una muestra no probabilística no alcanzan para generalizar a otras especies, tamaños, cámaras o puertos; las conclusiones valen para este montaje.

Hicimos la inferencia con una prueba de signos exacta sobre diferencias pareadas por réplica. No usamos pruebas paramétricas porque con $n = 6$ réplicas no se pueden verificar sus supuestos, y los conteos sobre 24 combinaciones son descriptivos. No hicimos la comprobación instrumental con barra patrón sumergida; la validación métrica es la comparación con las referencias de calibre. Tampoco cuantificamos en esta versión la repetibilidad del marcado, y por eso declaramos los objetivos 2 y 6 como parcialmente cumplidos.

\section{Recomendaciones}

Quien quiera repetir el procedimiento debería mantener fijos la cámara, la resolución, el enfoque, la iluminación, la posición de la pecera y la configuración del puerto durante todas las capturas, y conservar la imagen original, los puntos manuales, los archivos de calibración, la versión del programa y el registro de exclusiones. Recomendamos además registrar en cada sesión una medida física sumergida independiente, como una barra patrón, y repetir la evaluación con el pez en posiciones laterales del encuadre.

Para el uso práctico, la Tabla~\ref{tab:regla} resume qué hacer. La escala A es solo una línea de base y no corrige la refracción. B funciona si la posición del pez se expresa en el marco de su calibración y dentro del rango de distancias comprobado. C necesita conservar documentados los parámetros físicos del puerto, los índices de refracción y el sistema de coordenadas, y su salida solo es útil si la profundidad se conoce o se acota.

\section{Trabajos futuros}

\paragraph{Verificaciones inmediatas sobre este mismo conjunto.} Volver a marcar doce imágenes para medir la repetibilidad del operador, y hacer la comprobación con una barra patrón sumergida.

\paragraph{Líneas de continuación.} Quedan varios caminos abiertos. Detectar los puntos anatómicos de forma automática, con el mismo protocolo pareado, permitiría separar el error de localización del geométrico. Estimar la profundidad con una guía, una segunda vista o una restricción de postura reduciría el intervalo $\IL$ a un rango útil. También falta evaluar el método con peces vivos, agua turbia y posiciones fuera del centro del encuadre; incorporar al intervalo la incertidumbre de los parámetros del puerto y de los índices de refracción; y publicar un artículo con el conjunto de datos y el código.

\section{Impacto potencial y aporte}

El principal aporte de este trabajo es un procedimiento trazable para comparar conversiones métricas con una sola cámara en un montaje aire--vidrio--agua, y una medida concreta de lo que se puede afirmar sobre la longitud cuando no se conoce la profundidad. B y C fueron las rutas más exactas, con errores menores que el 1~\% cuando la profundidad es conocida, y \Aprime{} ofrece una alternativa más simple que mejora mucho a la escala fija en el vidrio.

Para la Ingeniería Informática, el trabajo deja cuatro productos: (1) un módulo de trazado refractivo de puerto plano, implementado a partir de las formulaciones de \textcite{Agrawal2012} y \textcite{Treibitz2012}, con pruebas unitarias de sanidad (con $\mu_g = \mu_w = 1$ y $t_g = 0$ reproduce la cámara pinhole) y evaluado contra referencias físicas; (2) una aplicación de escritorio de código abierto que conserva la trazabilidad de cada imagen (puntos, calibración, versión y advertencias) y exporta un registro reproducible; (3) un procedimiento experimental pareado que aísla el efecto del modelo geométrico y que sirve para evaluar, en el futuro, detectores automáticos de puntos; y (4) evidencia cuantitativa, para este montaje, sobre cuándo alcanza un marcador, cuándo conviene una calibración bajo el agua y cuándo hace falta el modelo físico.

Para la acuicultura, el resultado es una regla práctica, con números, para medir peces sin sacarlos del agua usando una cámara común a través del vidrio de un tanque. Llevarla a peces vivos o a instalaciones productivas sigue siendo una posibilidad futura, que requiere validarla primero en esas condiciones.
